\section{Introduction DRAFT}

Learning allows animals to use experience to predict future events and adapt their behavior accordingly. In classical conditioning, an initially neutral conditioned stimulus (CS) acquires behavioral significance through repeated pairing with a biologically salient unconditioned stimulus (US). The resulting conditioned response (CR) provides a quantitative readout of the learned relationship between cue and outcome. Because the timing, identity, and contingency of the stimuli can be experimentally controlled, classical conditioning remains a powerful framework for studying how nervous systems transform sensory experience into predictive behavior (Pavlov, 1927; Rescorla, 1988; Fanselow and Poulos, 2005; Gershman et al., 2010).

The temporal relationship between the CS and US is central to what the animal must learn. In delay conditioning, the CS overlaps with the US, allowing the association to be formed between events that are contiguous in time. In trace conditioning, by contrast, the CS ends before the US begins, leaving a stimulus-free interval that must be bridged behaviorally or neurally for learning to occur. Trace conditioning therefore provides a particularly useful assay for asking whether an animal can associate temporally separated events, rather than merely respond to overlapping stimuli. Establishing such learning in experimentally accessible preparations is important for linking behavioral evidence of temporal association to its underlying circuit mechanisms (Thompson and Steinmetz, 2009; Gilmartin and Helmstetter, 2010; Raybuck and Lattal, 2014; Kryukov, 2012).

Larval zebrafish offer a distinctive opportunity to study associative learning at brain-wide scale. Their small and optically transparent brains make them suitable for cellular-resolution functional imaging across large fractions of the nervous system, while their genetic and pharmacological accessibility enables mechanistic interrogation. At the same time, larvae produce quantifiable motor behaviors that can be monitored with high temporal resolution during controlled stimulation. These features make the model attractive for connecting learning behavior to distributed neural dynamics in an intact vertebrate brain (Ahrens et al., 2013; Portugues et al., 2014; Orger and de Polavieja, 2017; Vanwalleghem et al., 2018; Lovett-Barron, 2021).

Despite these advantages, classical conditioning paradigms in larval zebrafish still face methodological challenges. Existing studies have demonstrated associative learning in this model, but performance can be variable and may depend strongly on developmental stage, preparation, stimulus choice, and behavioral readout. Freely swimming assays preserve natural behavior but are less compatible with stable high-resolution imaging. Fully immobilized preparations improve optical stability but restrict behavioral output and can make learning difficult to measure. Head-fixed preparations, in which the head is stabilized while the tail remains free, offer a compromise between optical access and behavioral readout, but require assays that maintain sufficient baseline movement to detect learned changes in motor output (Aizenberg and Schuman, 2011; Valente et al., 2012; Harmon et al., 2017; Matsuda et al., 2017; Palumbo et al., 2020; Dempsey et al., 2022).

A further gap concerns trace conditioning. Delay conditioning has been demonstrated in larval zebrafish, but whether head-fixed larvae can acquire associations across a stimulus-free interval remains less clear. This question is important because trace conditioning extends the behavioral demands beyond simple temporal overlap between cue and outcome. Demonstrating trace learning in an imaging-compatible larval preparation would provide a foundation for future studies of how temporally separated events are represented and linked across the vertebrate brain. However, such a claim requires a robust assay with appropriate unpaired controls and careful separation of associative learning from changes in arousal, baseline activity, or non-associative responses to repeated stimulation (Aizenberg and Schuman, 2011; Harmon et al., 2017; Dempsey et al., 2022; Raybuck and Lattal, 2014).

Here, we developed a head-fixed larval zebrafish conditioning assay designed to preserve compatibility with future brain-wide imaging while retaining a measurable behavioral output. The assay pairs a visual CS with an aversive optovin-mediated US and quantifies behavior through tail-movement vigor. Optovin provides a temporally controlled stimulus that reliably evokes tail responses when activated by violet light, allowing repeated US delivery without mechanical contact. Because the learned response in this preparation is expressed as a reduction, rather than an increase, in tail movement, the protocol includes a priming and habituation structure intended to maintain an active baseline from which suppressive conditioned responses can be detected (Kokel et al., 2013; Aizenberg and Schuman, 2011; Harmon et al., 2017).

Using this framework, we tested whether larvae acquire both delay and trace associations. We show that head-fixed larvae develop learned suppression of tail-movement vigor after CS onset in delay conditioning and in a short-trace protocol with a 3-s stimulus-free interval. This suppression is absent in matched unpaired controls and diminishes during testing, supporting an associative interpretation. Under comparable conditions, conditioning was not detected at the population level with a 10-s trace interval, suggesting that longer temporal gaps impose a stronger constraint on this assay or preparation. Together, these findings establish predictive motor suppression as a behavioral readout of associative learning in head-fixed larval zebrafish and provide a foundation for future cellular-resolution studies of the neural mechanisms that support temporal association.
